\documentclass{article}

\usepackage[utf8]{inputenc}
\usepackage{helvet}
\usepackage{geometry}
\usepackage{xcolor}
\usepackage{eso-pic}
\usepackage{fancyhdr}
\usepackage{parskip}
\usepackage{microtype}
\usepackage[english, ukrainian]{babel}
\usepackage{url}
\usepackage{amsmath}
\usepackage{amssymb}
\usepackage{amsthm}
\usepackage{longtable}
\usepackage{booktabs}
\usepackage{hyperref}
\usepackage{titlesec}

\definecolor{nato}{HTML}{293757}
\definecolor{footergray}{gray}{0.65}

\geometry{a4paper,left=3cm,right=3cm,top=3cm,bottom=3cm}
\renewcommand{\familydefault}{\sfdefault}
\setcounter{secnumdepth}{3}

\DeclareFontFamily{T2A}{phv}{}
\DeclareFontShape{T2A}{phv}{m}{n}{<->ssub * cmss/m/n}{}
\DeclareFontShape{T2A}{phv}{bx}{n}{<->ssub * cmss/bx/n}{}
\DeclareFontShape{T2A}{phv}{m}{it}{<->ssub * cmss/m/it}{}
\DeclareFontShape{T2A}{phv}{bx}{it}{<->ssub * cmss/bx/it}{}

\titleformat{\section} {\color{nato}\Huge\bfseries}{\thesection.}{1em}{}
\titlespacing*{\section}{0pt}{30pt}{12pt}
\titleformat{\subsection} {\color{nato}\LARGE\bfseries}{\thesubsection.}{1em}{}
\titlespacing*{\subsection}{0pt}{20pt}{8pt}
\titleformat{\subsubsection} {\color{nato}\Large\bfseries}{\thesubsubsection.}{1em}{}
\titlespacing*{\subsubsection}{0pt}{10pt}{6pt}
\titleformat{\paragraph} {\color{nato}\large\bfseries}{\theparagraph.}{1em}{}
\titlespacing*{\paragraph}{0pt}{8pt}{4pt}


\fancyhf{}
\fancyhead[R]{\small\color{footergray} 27 червня 2026}
\fancyfoot[L]{\small\color{footergray} ERP/1. Операційна система}
\fancyfoot[R]{\small\color{footergray}\thepage}

\newcommand\HUGE[1]{\fontsize{40}{40}\selectfont #1}
\renewcommand{\headrulewidth}{0pt}
\renewcommand{\footrulewidth}{0.4pt}
\renewcommand{\footrule}{\color{footergray}\hrule width \textwidth height 0.4pt}

\usepackage{amsmath}
\usepackage{amssymb}
\usepackage{booktabs}
\usepackage{hyperref}
\usepackage{tikz}
\usetikzlibrary{shapes.geometric, arrows, arrows.meta, positioning, fit, backgrounds}

\let\originalfootnotesize\footnotesize
\let\oldverbatim\verbatim
\let\oldendverbatim\endverbatim
\renewenvironment{verbatim}{\originalfootnotesize\oldverbatim}{\oldendverbatim}

\title{\textbf{OS.1}: High-Assurance BEAM \\ on seL4 with Alpine VMM for drivers}
\author{Namdak Tonpa (maxim@synrc.com), Synrc Research}
\date{\today}

\hypersetup{
    colorlinks=true,
    linkcolor=blue,
    filecolor=magenta,
    urlcolor=cyan,
    pdftitle={OS.1 Hypervision},
    pdfpagemode=FullScreen,
}

\begin{document}

\begin{titlepage}
  \AddToShipoutPictureBG*{\AtPageLowerLeft{\color{nato}\rule{\paperwidth}{\paperheight}}}
  \color{white}\thispagestyle{empty}\vspace*{3cm}\centering
  {\Huge  \textbf{Zen Crypted Operating System} \par} \vspace{1.5cm}
  {\HUGE  \textbf{Технічне завдання} \par} \vspace{1.5cm}
  {\LARGE \textbf{на розробку та впровадження \\ операційної системи OS.1 \\ на гіпервізорі seL4} \par} \vspace{1.5cm}
  {\large На основі Технічних Вимог \par}
  {\large Згідно з Постановою КМУ № 205 від 21.02.2025 \par}
  {\large Формалізованих за стандартом ISO/IEC/IEEE 29148:2018 \par} \vfill
  {\large 2026 \copyright\ Криптографічні Телесистеми \par}
\end{titlepage}

\maketitle
\tableofcontents

\begin{abstract}
The Erlang/OTP BEAM runtime provides robust concurrency, fault isolation, and soft real-time behaviour that
make it attractive for high-availability systems, yet it traditionally depends on a large general-purpose
operating system whose trusted computing base (TCB) undermines strong security and certification arguments.
We present Tynse, a hybrid architecture that hosts an unmodified BEAM on the formally verified seL4
microkernel while retaining practical device support through a minimal Alpine Linux guest.

A purpose-built thin host (Tyn) implements only the small set of Linux-compatible system calls required
by a musl-linked OTP runtime and executes as an isolated seL4 protection domain under the Microkit framework.
Device drivers that would otherwise enlarge the TCB are confined to an Alpine Linux virtual machine managed
by a lightweight virtual-machine monitor; communication between the BEAM domain and the driver domain occurs
exclusively through capability-mediated shared-memory channels defined by the seL4 Device Driver
Framework (sDDF). The resulting system therefore combines the formal isolation guarantees of seL4,
the minimal attack surface of a BEAM-specific host, and the hardware coverage of an existing
Linux driver ecosystem.

We describe the static system architecture, the mapping of selected NIST SP~800-53 control families
onto seL4 capabilities and compile-time configuration, and a concrete bootstrap path on commodity
hardware (Raspberry Pi~4). The design demonstrates that a production-grade BEAM application can
operate with a dramatically reduced TCB while preserving the ability to utilise complex devices,
offering a practical route toward higher-assurance, certifiable distributed systems.
\end{abstract}

\section{Introduction}
\label{sec:intro}

Erlang/OTP and its BEAM virtual machine have long been valued for building highly concurrent,
fault-tolerant distributed systems. The language-level process model, supervision trees,
and soft real-time scheduling provide strong application-level isolation and recovery.
However, these benefits sit atop a conventional general-purpose operating system (typically Linux)
whose large trusted computing base (TCB) remains a significant obstacle to high-assurance and certification arguments.

Formal isolation at the kernel level has advanced substantially with the seL4 microkernel,
which offers a machine-checked proof of functional correctness and integrity. Yet running
a full Linux guest under seL4 merely relocates the TCB problem. Conversely, pure unikernel
approaches that eliminate the general-purpose OS often force substantial rewriting of
application runtimes or sacrifice device-driver coverage.

Tynse addresses this tension by combining three elements:

\begin{itemize}
  \item the formally verified seL4 microkernel and its Microkit framework for static, capability-based system construction;
  \item Tyn, a minimal host that supplies only the approximately fifty Linux-compatible system calls required by a musl-linked, unmodified OTP BEAM;
  \item a carefully confined Alpine Linux virtual machine that supplies real device drivers, communicating with the BEAM domain solely through seL4 Device Driver Framework (sDDF) shared-memory channels.
\end{itemize}

The resulting architecture keeps the TCB under the BEAM extremely small (seL4 + Tyn) while still permitting practical use of complex hardware. All isolation boundaries, scheduling parameters, and device ownership are fixed at image-construction time, yielding a static system amenable to analysis and aligned with selected families of NIST SP~800-53 controls.

This paper presents the architecture in detail, describes the principal protection domains and their interactions, and outlines the compile-time configuration model that governs resource allocation and scheduling.

\newpage
\section{Architecture}
\label{sec:architecture}

Tynse is organised as a static collection of seL4 protection domains (PDs) whose memory regions, communication channels, interrupt bindings and scheduling contexts are declared in a Microkit system description and materialised at build time. No protection domain or channel can be created dynamically at runtime.

\subsection{Components and Links}
\label{sec:components}

The principal components and the morphisms that connect them are summarised below.

\begin{center}
\noindent\makebox[\textwidth][c]{
\begin{tabular}{llp{10cm}}
\toprule
\textbf{Component} & \textbf{Type} & \textbf{Responsibility} \\
\midrule
seL4               & Microkernel          & Isolation, capabilities, MCS scheduling, IRQ routing \\
Microkit           & Framework            & Static system description $\to$ concrete capability layout \\
Tyn                & Protection Domain    & Syscall trap surface and minimal host for BEAM \\
BEAM / OTP         & Application runtime  & Erlang/Elixir processes, schedulers, supervision \\
VMM (libvmm)       & Protection Domain    & Creates and manages the Alpine Linux guest \\
Alpine Linux guest & Virtual Machine      & Real device drivers (network, block, \ldots) \\
UIO helper         & User-space (Linux)   & Bridges Linux drivers $\leftrightarrow$ sDDF rings \\
sDDF               & Protocol + libraries & Shared-memory rings and notifications \\
Console PD         & Protection Domain    & Early serial / console output \\
Monitor PD         & Protection Domain    & Metrics, health, audit events \\
Storage PD         & Protection Domain    & Block / SSD access (native or via Alpine) \\
Network PD         & Protection Domain    & Ethernet access (native or via Alpine) \\
\bottomrule
\end{tabular}
}
\end{center}

The dominant data and control paths are:

\begin{itemize}
  \item \textbf{Network}: NIC $\to$ Alpine driver $\to$ UIO helper $\to$ sDDF ring $\to$ Tyn $\to$ BEAM \texttt{gen\_tcp}/\texttt{Bandit};
  \item \textbf{Storage}: SSD/block device $\to$ Storage path $\to$ sDDF $\to$ Tyn $\to$ BEAM file or ETS/Mnesia;
  \item \textbf{Monitoring}: Tyn, VMM and driver PDs $\to$ notification + shared metrics region $\to$ Monitor PD;
  \item \textbf{Console}: UART/virtio-console owned exclusively by the Console PD exposed to Tyn.
\end{itemize}

All authority is expressed by seL4 capabilities granted at system construction. There is no ambient authority.

\newpage
\subsection{seL4 Microkit Capabilities}
\label{sec:microkit}

Microkit provides a deliberately constrained component model on top of seL4:

\begin{itemize}
  \item \textbf{Protection Domains (PDs)} --- isolated, single-threaded components with fixed address spaces (maximum 63 PDs in a system);
  \item \textbf{Memory Regions} --- explicitly mapped with read/write/execute rights into one or more PDs;
  \item \textbf{Channels} --- point-to-point connections supporting asynchronous notifications and synchronous protected procedure calls;
  \item \textbf{Virtual Machines} --- a PD may own at most one guest VM;
  \item \textbf{MCS scheduling contexts} --- each PD receives a budget and period; core affinity is declared statically.
\end{itemize}

The system description is processed at build time. The resulting image contains a fixed capability layout; the Microkit initialiser installs these capabilities before any PD begins execution. Consequently, the set of resources each component may access is known a priori and can be audited against security policy (including selected NIST SP~800-53 families such as AC, SC and CM).

Scheduling follows the seL4 Mixed-Criticality Scheduling (MCS) model. Typical priority ordering (highest to lowest) is:

\begin{enumerate}
  \item Console PD,
  \item Tyn PD (hosting BEAM),
  \item Monitor PD,
  \item VMM / Alpine driver VM,
  \item optional background domains.
\end{enumerate}

Core affinity and budget/period parameters are supplied via compile-time flags (\texttt{CORES}, \texttt{AFFINITY}, etc.)
and emitted into the system description.

\subsection{VMM Protection Domain}
\label{sec:vmm}

The VMM protection domain encapsulates the Linux driver environment so that complex device support does not enlarge the TCB under BEAM.

\subsubsection{libvmm}
\label{sec:libvmm}

libvmm is the lightweight virtual-machine monitor library used by Microkit. It runs as an ordinary protection domain and is responsible for:

\begin{itemize}
  \item allocating guest physical memory,
  \item establishing stage-2 translation,
  \item injecting interrupts,
  \item mediating hypercalls, and
  \item granting the guest only those device registers and IRQs that policy permits.
\end{itemize}

The VMM itself holds only the capabilities necessary to manage its single guest; it cannot map arbitrary host memory or interfere with Tyn or other native PDs.

\subsubsection{Alpine}
\label{sec:alpine}

Alpine Linux (musl-based, minimal) is chosen as the guest operating system for three reasons:
small footprint, binary compatibility with the musl-linked OTP that Tyn expects,
and straightforward packaging. A Buildroot- or minirootfs-derived image contains:

\begin{itemize}
  \item a stripped kernel configured for the target board and for UIO,
  \item a minimal user-space (busybox, openrc or equivalent),
  \item the UIO helper binary that maps sDDF rings and translates between Linux driver interfaces and the sDDF protocol.
\end{itemize}

Only the devices assigned to the guest appear in its device tree; all other hardware remains invisible.

\subsubsection{Tyn}
\label{sec:tyn}

Tyn executes as a sibling protection domain to the VMM, not inside the Linux guest.
It implements a deliberately narrow Linux-compatible system-call surface (approximately fifty calls)
that a musl-linked OTP binary expects: memory management (\texttt{mmap}, \texttt{mprotect}, \ldots),
threading and synchronisation (clone, futex family), time, basic file descriptors, and the in-memory
VFS that holds the OTP release.

Tyn never owns real device registers. All I/O is performed by mapping sDDF rings that are
filled by either native driver PDs or by the UIO helper running inside Alpine.
Consequently the trusted computing base that sits directly under BEAM consists solely of seL4 and Tyn.

\subsubsection{BEAM}
\label{sec:beam}

An unmodified (or minimally patched) musl-linked OTP release is embedded as a
cpio or tar archive inside the Tyn protection domain. At start-up Tyn locates
the archive, constructs the necessary environment, and transfers control to \texttt{beam.smp}.
From that point the ordinary OTP boot sequence proceeds: schedulers are started,
applications are loaded, and supervision trees become active. BEAM continues to
manage its own processes, message passing and distribution; seL4 only guarantees
the CPU budget and memory isolation of the enclosing Tyn domain.

\subsection{Monitor Protection Domain}
\label{sec:monitor}

The Monitor PD collects health and performance data from Tyn, the VMM and any native driver domains.
Communication occurs through a combination of seL4 notifications and a shared metrics memory
region whose layout is fixed at build time. The Monitor may:

\begin{itemize}
  \item expose counters and timestamps to an external observer,
  \item generate audit-style events (supporting NIST AU-family controls),
  \item detect liveness failures and signal higher-level recovery logic.
\end{itemize}

When the compile-time flag \texttt{MONITOR=0} is supplied the entire domain is omitted, further reducing image size and TCB.

\subsection{Console Protection Domain}
\label{sec:console}

The Console PD owns the platform UART (or a virtio-console device) exclusively.
It provides a narrow, character-oriented interface to Tyn for early boot messages,
panic output and optional interactive debugging.
Because the device registers reside only in the Console PD’s capability space,
neither Tyn nor the Alpine guest can access them directly.
The domain can be compiled to a write-only sink for production images.

\subsection{Storage Protection Domain}
\label{sec:storage}

Storage may be realised in two ways, selected at image-construction time:

\begin{itemize}
  \item \textbf{Native sDDF block driver} --- a dedicated PD owns the block device registers
        and IRQ and speaks the sDDF block protocol to Tyn;
  \item \textbf{Alpine-mediated} --- the block device is passed through to the Linux guest;
        the UIO helper translates Linux block I/O into sDDF rings.
\end{itemize}

In both cases Tyn presents a simple block or file-system abstraction to BEAM.
The choice is expressed by the \texttt{STORAGE} build flag and the corresponding
entries in the Microkit system description. Device ownership is exclusive:
a given controller is never mapped into more than one domain.

\subsection{Network Protection Domain}
\label{sec:network}

Analogously to storage, networking is provided either by a native sDDF Ethernet driver PD or by the Alpine guest.
The data path is identical in structure:

\begin{center}
NIC $\to$ (native PD or Alpine + UIO) $\to$ sDDF ring $\to$ Tyn $\to$ BEAM socket layer.
\end{center}

Compile-time selection (\texttt{NET=1}) determines which implementation is linked into the final image.
Core affinity for the network domain can be set independently so that packet processing does
not perturb BEAM’s soft real-time behaviour.

\newpage
\subsection{NIST SP 800-53 Control Mapping}
\label{sec:nist}

Tynse’s static, capability-based design maps naturally onto several families of NIST SP~800-53 controls. The table below summarises the primary alignments. This is a design-time correspondence, not a formal certification claim.

\begin{center}
\noindent\makebox[\textwidth][c]{
\begin{tabular}{p{6cm}p{3cm}p{8cm}}
\toprule
\textbf{Family} & \textbf{Example Controls} & \textbf{How Tynse addresses them} \\
\midrule
\textbf{AC} Access Control
  & AC-3, AC-4, AC-6
  & seL4 capabilities and Microkit channels enforce least privilege and information-flow control; no ambient authority exists at runtime. \\
\addlinespace
\textbf{AU} Audit \& Accountability
  & AU-2, AU-12
  & Monitor PD collects events and metrics over explicit notification channels; audit data can be exported or retained in a dedicated region. \\
\addlinespace
\textbf{CM} Configuration Management
  & CM-2, CM-6, CM-7
  & All isolation boundaries, device ownership and feature sets are fixed at image-construction time; unused components are compiled out. \\
\addlinespace
\textbf{SC} System \& Communications Protection
  & SC-7, SC-39
  & Strong spatial isolation between PDs; boundary protection is realised by capability checks rather than software firewalls. \\
\addlinespace
\textbf{SI} System \& Information Integrity
  & SI-7, SI-16
  & Static, reproducible images; capability-based memory protection; optional measured-boot hooks. \\
\addlinespace
\textbf{CP} Contingency Planning
  & CP-9, CP-10
  & BEAM supervision trees provide deterministic process-level recovery inside the Tyn domain; seL4 contains failures to the originating PD. \\
\addlinespace
\textbf{IA} Identification \& Authentication
  & IA-2, IA-5
  & Optional; can be realised as a dedicated protection domain that authenticates external principals before granting channels. \\
\addlinespace
\textbf{SA} System \& Services Acquisition
  & SA-11, SA-15
  & Small TCB, open design, reproducible builds, and clear separation of upstream components (seL4, Microkit, Tyn, Alpine). \\
\bottomrule
\end{tabular}
}
\end{center}

The compile-time flags (\texttt{SECURITY=nist53}, \texttt{MONITOR=1}, \texttt{CORES}, \texttt{AFFINITY}, etc.) directly influence the strength of the above controls by determining which domains are present, which devices are accessible, and how CPU time is partitioned.

\newpage
\section{Conclusion}
\label{sec:conclusion}

Synrc Hypervision demonstrates that a production-grade, unmodified BEAM runtime can be hosted on the formally verified seL4
microkernel with a dramatically reduced trusted computing base while still retaining practical access to complex
devices. The architecture achieves this by:

\begin{itemize}
  \item confining the BEAM to a minimal syscall-trap host (Tyn) that runs as an ordinary seL4 protection domain;
  \item isolating device drivers inside a capability-restricted Alpine Linux virtual machine (or native sDDF PDs);
  \item mediating all inter-component communication through static, capability-controlled sDDF channels;
  \item fixing scheduling parameters, core affinity and device ownership at image-construction time.
\end{itemize}

The resulting system is static, analysable, and naturally aligned with key NIST SP~800-53 control families.
Future work includes completing native sDDF drivers for the Raspberry Pi platform, strengthening measurement
and attestation support, and evaluating end-to-end performance and fault-isolation properties under realistic OTP workloads.

\end{document}
